\subsection{Presupuesto conjunto de las dos arquitecturas}
\label{sec:source-budget}
Sea \(q\) el número total de bloques liberados y sea \(Z_i\) la vista para la
que se compromete el servicio, incluidos el estado inicial y el historial
pertinente, las alarmas y la telemetría accesible. La selección y el
acondicionamiento pueden depender del régimen físico. Si, en cada paso,
el modelo respalda \(H_{\min}(R_i\mid Z_i)\ge k_i\) en el modelo ideal
de selector independiente, y la ley real de selector, dato y vista se aparta
del modelo a lo sumo \(\eta_i\), un presupuesto suficiente es
\[
 \epsilon_{\rm flujo}\ \le\ \min\left(1,
 \sum_{i=1}^{q}\bigl[2^{(128-k_i)/2-1}+2\eta_i\bigr]\right).
\]
La vista de cada término cubre las condiciones de composición de la
proposición de secuencia. El operador suma también los bloques utilizados
para resiembras si reivindica la calidad de dicho material. Se especifican por
separado el error de la fuente, la distancia estadística del flujo y las
ventajas computacionales DRG.3.

Para \(q=2^{32}\) bloques, objetivo \(2^{-64}\) y selector ideal, despejar
\(q\,2^{(128-k)/2-1}\le2^{-64}\) da \(k\ge318\). Es el mínimo
entero de \emph{esta cota suficiente}. Con \(k=324\) el término de extracción
suma \(2^{-67}\). Por ejemplo, si \(\eta_i\le2^{-98}\) para cada paso,
el coste adicional es \(2^{-65}\) y el total cabe en \(2^{-64}\).
Con \(k=268\) el mismo cálculo sólo da \(2^{-39}\) por la extracción en
todo el horizonte. El control por matriz de §338(c) y el compromiso ITS del
flujo tienen cuantificadores distintos. Los certificados enteros
\proofat{PTG.SourceBudget.exactThreshold}{exactThreshold}, \proofat{PTG.SourceBudget.idealMargin}{idealMargin} e \proofat{PTG.SourceBudget.imperfectSelectorBudget}{imperfectSelectorBudget}
se encuentran en \proofmodule{SourceBudget}.

\paragraph{Dos controles del modelo.}
Primero, considérese una palabra formada por 126 bits uniformes independientes
seguidos de copias de los dos primeros. Cada bit está equilibrado y el bloque
tiene exactamente 126 bits de minentropía, tasa superior a 0.98.
Su ley uniforme sobre un subespacio ocupa un cuarto del alfabeto y su distancia
a \(U_{128}\) es \(3/4\). Ambas biyecciones, fijado el estado, conservan esa
distancia. Esto explica por qué el valor de PTG.2.2 aislado no determina
un error de servicio. Segundo, elíjase un selector uniforme \(U\) y,
condicionado a él, \(R\) uniforme en \(\ker T_U\).
Por rango y nulidad, cada núcleo tiene al menos \(2^{n-128}\) elementos.
Para \(n\ge452\) se cumplen hasta 324 bits de minentropía condicionada
a \(U\), mientras que \(T_UR=0\) en todas las realizaciones. La prueba de
extracción necesita justificar la dependencia conjunta. \proofat{PTG.SourceBudget.kernelSize}{kernelSize}
y \proofat{PTG.SourceBudget.dependentSourceConstant}{dependentSourceConstant} verifican en Lean un ejemplo exhaustivo
de una fila y tres columnas. El argumento para \(128\times n\) usa rango y
nulidad. Un estudio experimental para la fuente modelará por separado las
adquisiciones de selector y datos, su separación temporal y el estado del
control. Una estimación de la minentropía de cada canal por separado
resulta insuficiente \cite[§§335, 338 y 951--955]{ais}\cite{lhl}.

\paragraph{Alternativa de interfaz por lotes.}
Se puede investigar una extracción ancha de \(128b\) bits con un selector
renovado por lote, seguida de \(b\) llamadas completas con subbloques
disjuntos de 128 bits. La inversa secuencial conserva la calidad conjunta
del lote. El presupuesto de extracción sustituye 128 por \(128b\) y el
número de renovaciones por el número de lotes, incluyendo el último lote
parcial en la contabilidad. Para \(b\in\{1,4,16\}\) las anchuras son 128,
512 y 2048. Se condiciona el cambio a un modelo de ventana más larga, a la
protección del búfer y a la aceptación de la unidad escogida como número
intermedio según §§336,338,385--390. La operación de referencia en los dos
expedientes es \(b=1\).
